\documentclass[12pt,a4paper]{article}
\usepackage[utf8]{inputenc}
\usepackage[T1]{fontenc}
\usepackage[portuguese]{babel}
\usepackage{graphicx}
\usepackage{hyperref}
\usepackage{geometry}
\geometry{margin=2.5cm}

\title{Sistema de Gestão de Tarefas com Arquitetura de Microsserviços}
\author{Nome do Aluno / Grupo}
\date{\today}

\begin{document}
\maketitle

\section{Introdução}
Este trabalho apresenta o desenvolvimento de um sistema de gestão de tarefas baseado em arquitetura de microsserviços. O sistema foi dividido em serviços independentes para autenticação, gestão de utilizadores, gestão de tarefas e orquestração dos pedidos.

\section{Objetivos}
O objetivo principal é implementar uma aplicação REST com microsserviços, Docker e Kubernetes local através do Minikube. O sistema permite registar utilizadores, autenticar com JWT e gerir tarefas associadas ao utilizador autenticado.

\section{Arquitetura do Sistema}
A arquitetura adotada segue o modelo de orquestração. Todos os pedidos externos são enviados para o Orchestrator Service, que encaminha os pedidos para o Auth Service, User Service ou Task Service.

\begin{verbatim}
Cliente/Postman
      |
      v
Orchestrator Service
      |
      |--> Auth Service
      |--> User Service
      |--> Task Service
              |
              v
            MongoDB
\end{verbatim}

\section{Microsserviços Implementados}
\subsection{Orchestrator Service}
Responsável por receber todos os pedidos externos e reenviá-los para os serviços apropriados.

\subsection{Auth Service}
Responsável pela autenticação dos utilizadores e geração de tokens JWT.

\subsection{User Service}
Responsável pela criação, consulta, atualização e remoção de utilizadores.

\subsection{Task Service}
Responsável pela criação, consulta, atualização e remoção de tarefas.

\section{Base de Dados}
Foi utilizado MongoDB como base de dados não relacional. O serviço de utilizadores utiliza a base de dados \texttt{usersdb} e o serviço de tarefas utiliza a base de dados \texttt{tasksdb}.

\section{Docker}
Cada microsserviço possui o seu próprio \texttt{Dockerfile}, permitindo que cada serviço seja executado de forma independente num container.

\section{Kubernetes e Minikube}
Os serviços foram integrados num ambiente Kubernetes local usando Minikube. Foram criados ficheiros YAML para deployments e services de cada microsserviço.

\section{Testes}
Os testes foram realizados através do Postman/curl. Foram testadas as seguintes funcionalidades: registo, login, criação de tarefas, listagem de tarefas, atualização e eliminação de tarefas.

\section{Conclusão}
O projeto demonstrou a implementação de um sistema distribuído com microsserviços, comunicação REST, autenticação com JWT, Docker e Kubernetes local com Minikube.

\end{document}
